\chapter{RL Circuits}
\label{ch:rl}
An RL circuit is the current-domain sibling of the RC circuit: it is the same
\emph{first-order system} of \cref{ch:linsys}, but now the natural state is the
inductor current \(i_L\), which holds the stored magnetic energy.

\section{Where RL Circuits Appear}
Inductors most often appear \emph{with} capacitors, in the resonant circuits of
\cref{ch:rlc,ch:rlcpar}---so it is fair to ask where an inductor works against a
resistance alone. Several places, and they matter:
\begin{description}
\item[RF chokes] a coil in a bias or supply lead, there to pass DC while blocking
RF. What matters is the choke's inductance against the resistance of the path it
feeds, and \(\tau=L/R\) sets how well it isolates.
\item[Relay and solenoid coils] a coil switched by a transistor is an inductance in
series with its own winding resistance, nothing more. Its \(L/R\) governs how fast
it pulls in---and its stored energy is why the switch needs a flyback diode.
\item[Transformer and motor windings] modeled as leakage inductance in series with
winding resistance, which is what limits how fast the current can change.
\item[Loading coils and conductor loss] the coil in a shortened antenna is an
inductance whose unavoidable series resistance sets its \(Q\), directly limiting
efficiency and bandwidth (\cref{ch:lines}).
\item[Every real inductor] a coil \emph{is} an RL circuit by itself, since its
winding resistance sits in series with its inductance; its self-capacitance only
matters near self-resonance.
\end{description}
The RL circuit therefore earns its own chapter both because these cases are real and
because the RC/RL duality it exposes is the fastest way to remember either one.

\section{The Circuit}
A source \(u(t)\) drives a resistor \(R\) in series with an inductor \(L\)
(\cref{fig:rl-schematic}). The state is the inductor current \(i\), but the
\emph{output} is a \textbf{voltage}---the one across the resistor:
\[
v_o=v_R=R\,i .
\]
Since \(v_R\) is simply \(R\) times the state, this output is a scaled picture of
the inductor current, which is why it is called low-pass ``in the current sense.''
Taking \(v_o\) across the inductor instead gives a high-pass response.

\begin{figure}[htbp]
\centering
\begin{circuitikz}[scale=1]
\draw (0,0) to[sV,l=\(u(t)\)] (0,2)
      to[L,l=\(L\),i>=\(i(t)\)] (3,2)
      to[R,l=\(R\),v=\(v_R\)] (3,0) -- (0,0);
\draw (3,2) to[short,-o] (4,2) node[right]{$+$};
\draw (3,0) to[short,-o] (4,0) node[right]{$-$};
\node at (4.35,1) {\(v_o\)};
\end{circuitikz}
\caption{Series RL circuit. The state variable is the inductor current \(i(t)\);
the output is the \emph{voltage} across the resistor, \(v_o=v_R=Ri\), which is
low-pass in the current sense.}
\label{fig:rl-schematic}
\end{figure}

\section{From Schematic to Differential Equation}
KVL gives the source as the sum of the inductor and resistor voltages, and the
inductor law relates its voltage to its rate of current change:
\[
u(t)=L\frac{\dd i}{\dd t}+R\,i(t)
\qquad\Longleftrightarrow\qquad
\boxed{\,\frac{L}{R}\frac{\dd i}{\dd t}+i=\frac{1}{R}\,u(t).\,}
\]

\section{State-Space Form and the Generic Match}
With the single state \(x=i\), and remembering that the output is the resistor
voltage,
\[
\dot{x}=-\frac{R}{L}\,x+\frac{1}{L}\,u,
\qquad
y=v_o=v_R=R\,x .
\]
The output equation is what distinguishes this from a bare current model: \(y\) is
the resistor voltage, so \(C_{\text{out}}=R\) and \(D=0\). Transforming the state
equation gives the transfer function to the current, and multiplying by \(R\) gives
the one to the output:
\[
\frac{I(s)}{U(s)}=\frac{1/L}{s+R/L},
\qquad
G(s)=\frac{V_o(s)}{U(s)}=R\cdot\frac{1/L}{s+R/L}=\frac{1}{1+(L/R)s},
\]
a unity-DC-gain first-order low-pass---the same form as the RC circuit's
\(G_C\). Comparing the state equation with \(\tau\dot y+y=Ku\) from
\cref{ch:linsys} (taking \(y\) as the \emph{current}, whose DC gain is \(1/R\)),
\[
\tau=\frac{L}{R},
\qquad
K=\frac{1}{R},
\qquad
s_{\text{pole}}=-\frac{R}{L}=-\frac{1}{\tau}.
\]
It is the same first-order system as the RC circuit, with the time constant
\(L/R\) in place of \(RC\). Note the duality: RC stores energy in a voltage and
its time constant grows with \(R\); RL stores energy in a current and its time
constant \emph{shrinks} with \(R\).

\section{Transient Response}
For a DC step of height \(E\), and for decay from an initial current \(I_0\),
\[
i(t)=\frac{E}{R}\left(1-\ee^{-Rt/L}\right),
\qquad
i(t)=I_0\,\ee^{-Rt/L}.
\]
This is again the first-order response of \cref{ch:linsys}, read as an inductor
current (\cref{fig:rl-transient}): after one time constant \(\tau=L/R\) the
current has reached \(63.2\%\) of its final value. Inductor current cannot jump
under finite voltage---the dual of the capacitor's continuous voltage.

\begin{figure}[htbp]
\centering
\includegraphics[width=0.82\textwidth]{rl_transient}
\caption{RL inductor-current transient. Rise and decay share the time constant
\(\tau=L/R\); one \(\tau\) covers \(63.2\%\) of the change. The same shape as the
RC voltage transient---the RC/RL duality.}
\label{fig:rl-transient}
\end{figure}

\section{Frequency Response and the Root Locus}
The inductor impedance is \(Z_L=\jj\omega L\), which grows with frequency, so an
inductor opposes rapid current changes and looks like a short at DC. Output across
\(R\) is low-pass (\cref{fig:bode1}); across \(L\) it is high-pass, identical in
shape to the RC high-pass of \cref{fig:rc-highpass}, with corner
\(\omega_c=R/L\). The single pole \(-R/L\) rides the negative real axis with the
same geometry as the RC pole (\cref{fig:rc-rootlocus}); as \(R\)
increases the pole moves \emph{left} (faster), the mirror of the RC case.

\begin{workedbox}[: sizing a choke]
A bias line feeding a \SI{50}{\ohm} path must pass DC while presenting at least
\SI{500}{\ohm} of reactance at \SI{14}{\mega\hertz}. From \(X_L=\omega L\),
\[
L\ge\frac{X_L}{\omega}=\frac{500}{2\pi(14\times10^6)}\approx\SI{5.7}{\micro\henry},
\]
so a standard \SI{10}{\micro\henry} choke (\(X_L\approx\SI{880}{\ohm}\)) is
comfortable. Read the same part as a first-order system and the choice is the
corner frequency: \(\tau=L/R=\SI{0.2}{\micro\second}\) puts \(f_c\approx
\SI{796}{\kilo\hertz}\), and \SI{14}{\mega\hertz} is \(1.25\) decades past it,
hence about \SI{25}{\decibel} of attenuation. Two vocabularies, one pole.
\emph{But check \cref{sec:selfresonance} before trusting it:} a
\SI{10}{\micro\henry} coil with \SI{2}{\pico\farad} of winding capacitance
self-resonates near \SI{36}{\mega\hertz}, above which it stops being an inductor
at all.
\end{workedbox}

\begin{controlsbox}
\textbf{Connection.} RC and RL are the two first-order lags of circuit theory.
Same pole structure, same Bode shape, same step response---only the state
variable (voltage vs current) and the direction of the \(R\)-dependence differ.
Learn one and you have both.
\end{controlsbox}

\section{Ferrites and Common-Mode Current}
\label{sec:commonmode}
The most common inductor in a real station is one nobody drew: a ferrite bead or
toroid clamped over a cable. Understanding it needs one new idea---not about
inductors, but about which current the inductor is in series with.

\subsection{Two Currents on the Same Cable}
Take any cable with two or more conductors. The currents in them can be decomposed
into two independent parts:
\begin{description}
\item[Differential mode] equal and opposite in the two conductors. This is the
wanted signal: current out along one conductor and back along the other, with the
two magnetic fields cancelling outside the cable.
\item[Common mode] equal and in the \emph{same} direction in every conductor,
returning by some other path entirely---through the ground, the mains wiring, or the
operator. Its fields do \emph{not} cancel, so a common-mode current turns the cable
into an antenna.
\end{description}
That is the whole diagnosis of a great deal of interference. A feed line radiating,
a microphone lead picking up RF, a computer cable polluting the band: in each case
the differential signal is fine and a common-mode current is doing the damage.

\subsection{Why a Ferrite Choke Is Selective}
Now the circuit question: how do you add impedance to one mode and not the other?
Thread the whole cable through a ferrite core. The core links whatever net flux the
conductors produce, and that is precisely the distinction:
\begin{itemize}
\item For the \emph{differential} current the two conductors' fields cancel, so
there is no net flux, so the core adds \emph{no} inductance. The wanted signal
passes untouched.
\item For the \emph{common-mode} current all the conductors push flux the same way,
so the core sees the full sum and presents a series impedance \(\jj\omega L\) to it.
\end{itemize}
So a ferrite choke is an inductance that exists only for the current you want to
suppress---not a filter in frequency, but a filter in \emph{mode}. It is the RL
circuit of this chapter with the common-mode path's resistance playing the role of
\(R\), and \(\tau=L/R\) again setting where the suppression begins.

Two practical consequences follow from the formulas already in hand. More turns
through the core raises \(L\) (roughly as the square of the turn count, since both
the flux and the linking scale), which lowers the frequency at which the choke
starts working. And a ferrite's \textbf{mix}---its material composition---determines
the frequency range over which it is lossy rather than merely reactive; a lossy mix
is preferable here, because dissipating the common-mode energy is better than
storing and re-radiating it.

\begin{physicalbox}
\textbf{Why a toroid, and not a rod.} A closed core keeps its flux inside the
magnetic circuit, so almost all of it links the winding: that gives a large
inductance per turn and, just as usefully, little external field to couple into
neighboring circuits. An open rod leaks flux through the air, so it needs far more
turns for the same inductance and it broadcasts to its surroundings. The same
containment argument is why toroidal cores dominate in RF work generally---the
inductance is high, the coupling to everything else is low, and the mix can be
chosen per band.
\end{physicalbox}

\section{Energy Interpretation}
The inductor stores \(E_L=\tfrac12 Li^2\) in its magnetic field. When current is
interrupted, that energy must go somewhere.

\begin{energybox}
Because \(E_L=\tfrac12Li^2\) cannot vanish instantly, breaking an inductive current
forces the inductor to produce whatever voltage is needed to keep the current
flowing---which is why an opening switch arcs, and why flyback diodes and snubbers
exist. The dual statement for a capacitor is that shorting it produces whatever
\emph{current} is needed to keep its voltage continuous.
\end{energybox}

\begin{exambox}
Remember the natural state: capacitor voltage and inductor current. That one
sentence tells you which quantity is continuous and which differential equation
to write.
\end{exambox}

\section{Worked Example}
A \(\SI{10}{\micro\henry}\) inductor is in series with a \(\SI{50}{\ohm}\)
resistor. The time constant and pole are
\[
\tau=\frac{L}{R}=\frac{10\times10^{-6}}{50}=\SI{2.0e-7}{\second},
\qquad
s_{\text{pole}}=-\frac{R}{L}=-\SI{5e6}{\per\second},
\]
so the corner frequency is \(f_c=1/(2\pi\tau)\approx\SI{796}{\kilo\hertz}\). At
\SI{14}{\mega\hertz}---well above the corner---the reactance
\(X_L=2\pi(14\times10^6)(10\times10^{-6})\approx\SI{880}{\ohm}\) dominates the
\SI{50}{\ohm} resistor, so the inductor's impedance sets the behavior.

\Cref{fig:rl-worked} plots this specific circuit. On the magnitude curve the
\SI{14}{\mega\hertz} point sits about \SI{25}{\decibel} down---consistent with being
roughly \(1.25\) decades past the corner at \(-\SI{20}{\decibel}\) per decade---which
is exactly the choke behavior this combination is used for: it passes DC and audio
while attenuating HF.

\begin{figure}[htbp]
\centering
\includegraphics[width=\textwidth]{rl_worked}
\caption{The worked example, \(L=\SI{10}{\micro\henry}\) with
\(R=\SI{50}{\ohm}\). \emph{Left:} the Bode plot of \(v_R/u\) on a real frequency
axis, with the \SI{796}{\kilo\hertz} corner and the \SI{14}{\mega\hertz} operating
point marked. \emph{Right:} the root locus as \(R\) is swept with \(L\) fixed. The
pole \(s=-R/L\) moves \emph{left} (faster) as \(R\) grows---the mirror of the RC
case, where larger \(R\) moved the pole right.}
\label{fig:rl-worked}
\end{figure}

\begin{physicalbox}
\textbf{How does adjusting \(R\) change the dynamics?} As with the RC circuit, the
single real pole can only slide along the real axis---no value of \(R\) produces
ringing. But the \emph{direction} is reversed: here \(\tau=L/R\), so increasing
\(R\) \emph{shortens} the time constant and moves the pole left (faster, higher
cutoff), whereas in the RC circuit increasing \(R\) lengthened \(\tau\) and moved
the pole right. This is the RC/RL duality in one sentence, and it is the detail most
often misremembered: for a choke, \emph{lowering} the circuit resistance is what
makes the isolation better at low frequencies.
\end{physicalbox}
